\section{Related Work}
\label{sec:related_work}

\subsection{Molecular Property Prediction}
Molecular property prediction has been standardized through benchmark suites such as MoleculeNet~\cite{wu2018moleculenet}, which provides curated classification and regression tasks spanning physicochemical, biophysical, and physiological domains. State-of-the-art approaches include message-passing neural networks (MPNNs)~\cite{gilmer2017neural}, graph transformer architectures~\cite{ying2021transformers}, and descriptor-based models using extended-connectivity fingerprints~\cite{rogers2010extended}. Scaffold splitting~\cite{bemis1996properties} is widely adopted to evaluate generalization under structural novelty. \textbf{Distinction:} Our work is not a model development study. We do not propose a new prediction architecture. Rather, we use an existing LLM in a zero-shot setting to investigate how different types of textual augmentation affect prediction behavior.

\subsection{LLMs in Molecular Science}
Large language models have been increasingly applied to chemical tasks, including reaction prediction~\cite{schwaller2019molecular}, retrosynthesis~\cite{liu2024retrosynthesis}, and molecular property estimation~\cite{guo2024llm, ross2022large}. Text-based molecular representations, where SMILES strings are accompanied by natural language descriptions, have shown promise for zero-shot and few-shot inference~\cite{edwards2022text2mol}. Recent studies have explored whether LLMs encode implicit chemical knowledge through their pre-training corpora~\cite{white2023assessment}. \textbf{Distinction:} Our work is not a prompt engineering study. We do not optimize prompts for maximum prediction accuracy. Instead, we use a fixed prompt template to systematically vary the semantic content of augmentation text while controlling for confounding factors.

\subsection{Hallucination in Scientific LLM Applications}
LLM hallucination---the generation of fluent but factually unsupported text---has been extensively studied in general NLP contexts~\cite{ji2023survey, huang2023survey}. In molecular and scientific applications, hallucination poses particular risks because factual errors in structural or mechanistic claims may not be easily detected by non-expert users~\cite{ramos2024review}. Recent work by Hao et al.~\cite{hao2024hallucination} reported the counterintuitive finding that hallucinated descriptions can improve molecular property prediction, though without characterizing which aspects of the hallucinated text drive this effect. Yuan et al.~\cite{yuan2025taxonomy} proposed a taxonomy for molecular hallucinations but did not conduct controlled ablation experiments isolating the contribution of each category. \textbf{Distinction:} Our work is not a hallucination detection or mitigation study. We treat hallucinations as structured semantic perturbations and study their differential effects on prediction behavior through controlled ablation, incorporating semantic and stylistic controls absent from prior evaluations.
